\chapter[CONCLUSÕES FINAIS E TRABALHOS FUTUROS]{CONCLUSÕES FINAIS E TRABALHOS FUTUROS}
\label{cha:conclusao}

% -------------------------------------------------------------------------
\section{Conclusões}
\label{sec:conclusoes}

Este trabalho investigou a viabilidade de utilizar o Kubernetes — especificamente o padrão de Operator e o mecanismo de Definições de Recursos Customizados (CRDs) — como plataforma de orquestração para aplicações de controle distribuído modeladas segundo os princípios do IEC~61499. A prova de conceito foi realizada sobre uma implementação do Processo Tennessee Eastman (TEP) em Rust, exposta via gRPC, com supervisão Kubernetes em Go/Kubebuilder e interface de operação em FastAPI/WebSocket.

Os principais resultados obtidos permitem afirmar:

\begin{enumerate}

  \item \textbf{Viabilidade técnica da orquestração Kubernetes para controle industrial.}
  O Operator Kubernetes demonstrou capacidade de observar continuamente o estado de uma planta industrial simulada, avaliar variáveis de processo contra uma política declarativa e aplicar ações corretivas via gRPC — tudo de forma autônoma e sem intervenção humana. A fase \texttt{Stable} foi mantida por horizonte indeterminado (validado em mais de 30.000~h simuladas), e a latência de reconciliação (escala de segundos) é compatível com a dinâmica lenta do TEP.

  \item \textbf{O CRD \texttt{PLCMachine} como equivalente declarativo do modelo de dispositivo IEC~61499.}
  A separação entre \textit{spec} (política desejada) e \textit{status} (estado observado) do CRD espelha a separação entre modelo de aplicação e estado de execução do IEC~61499. O Operator assume o papel do runtime de implantação e reconfiguração do padrão, tornando o Kubernetes uma infraestrutura viável para o gerenciamento de aplicações de controle distribuído.

  \item \textbf{A separação entre modelo de planta e lógica de controle habilita gestão externa.}
  A refatoração do Experimento~11, que extraiu os controladores para a trait \texttt{Controller} e o \texttt{ControllerBank}, foi o passo arquitetural que tornou possível a gestão dos controladores pelo Operator. Esse resultado reforça a importância da separação de responsabilidades em sistemas de controle distribuído — princípio central do IEC~61499.

  \item \textbf{Controladores P são insuficientes para rejeitar IDV(1) e IDV(2).}
  Os Experimentos~15 e 16 demonstraram que os três controladores P do baseline, embora suficientes para estabilidade nominal, não possuem autoridade de controle suficiente para rejeitar distúrbios composicionais. IDV(1) colapsa a planta em $\approx$2,5~h por mecanismo cinético; IDV(2) causa colapso lento ($\approx$30~h) por acúmulo de inerte. Ambos os casos demandam lógica supervisória de nível mais alto — exatamente o papel do Operator.

  \item \textbf{IDV(3) é numericamente ineficaz nesta implementação.}
  O Experimento~17 confirmou que a perturbação de entalpia do D~feed é diluída pelo reciclo antes de atingir o reator, não gerando resposta mensurável. Esse resultado é consistente com a análise do código FORTRAN original e delimita o espaço de distúrbios relevantes para validação do supervisor a IDV(1) e IDV(2).

  \item \textbf{A arquitetura é extensível e substituível.}
  A comunicação baseada em contratos gRPC tipados (\texttt{.proto}) garante que qualquer componente pode ser substituído — incluindo o simulador Rust por dados de uma planta física real — sem alteração dos demais. Isso valida a arquitetura como protótipo realista para cenários industriais.

\end{enumerate}

Em síntese, o trabalho demonstra que o Kubernetes, com seu modelo declarativo e o padrão de Operator, constitui uma plataforma tecnicamente viável para orquestração de aplicações de controle distribuído compatíveis com os princípios do IEC~61499 — respondendo afirmativamente à questão de pesquisa proposta.

% -------------------------------------------------------------------------
\section{Trabalhos Futuros}
\label{sec:futuros}

A infraestrutura desenvolvida estabelece uma base sobre a qual as seguintes extensões são diretamente aplicáveis:

\begin{enumerate}

  \item \textbf{Avaliação automática de qualidade de malhas — Índice de Preditividade.}
  A próxima evolução natural do Operator é incorporar o cálculo do Índice de Preditividade (PI) proposto pelo CERN \cite{CERN_WEAPL02} sobre as séries temporais das XMEAS. O Operator passaria de observador puro para avaliador de desempenho: identificaria automaticamente quais malhas estão oscilando ou com resposta degradada e registraria esse diagnóstico no \textit{status} da \texttt{PLCMachine}. O gêmeo digital, com controle total sobre os distúrbios, é o ambiente ideal para gerar os dados de operação degradada necessários para calibrar e validar o índice.

  \item \textbf{Intervenção ativa via \texttt{UpdateController}.}
  Com o diagnóstico de qualidade estabelecido, o Operator pode ser estendido para agir: detectar tendência de pressão crescente (IDV(1), IDV(2)) e ajustar setpoints ou ganhos dos controladores antes que o erro acumulado cause intertravamento. A infra\-estrutura gRPC e o método \texttt{UpdateController} já estão implementados — falta a lógica de decisão.

  \item \textbf{Integração OPC UA para interoperabilidade.}
  O OPC UA (\textit{Unified Architecture}), padronizado na IEC~62541, é o protocolo de interoperabilidade industrial amplamente adotado em ambientes heterogêneos. Desenvolver um servidor OPC UA sobre a planta simulada permitiria conectar controladores reais, ferramentas SCADA e sistemas de terceiros ao gêmeo digital sem alteração da camada de simulação — transformando o laboratório em uma plataforma de integração e teste de sistemas industriais reais.

  \item \textbf{Exploração dos frameworks do CERN — EPICS e TANGO Controls.}
  O CERN mantém dois ecossistemas de controle abertos amplamente usados em física experimental: EPICS (\textit{Experimental Physics and Industrial Control System}) e TANGO Controls. Ambos modelam dispositivos com variáveis de processo (PVs), comandos e estados — estrutura diretamente mapeável às XMEAS e XMV do TEP. Integrar o gêmeo digital a um desses frameworks abriria acesso a décadas de ferramentas de diagnóstico, arquivamento e supervisão já consolidadas na comunidade científica.

  \item \textbf{Kubernetes na borda com KubeEdge ou k3s.}
  Para uso em cenários industriais reais, o Operator deve executar próximo ao hardware de campo. Distribuições leves como k3s ou o framework KubeEdge são adequados para PLCs e gateways industriais, permitindo que a mesma lógica supervisória desenvolvida aqui seja implantada em produção sem reescrita.

  \item \textbf{Modelagem completa do trocador de calor do reator.}
  O Experimento~14 identificou que IDV(4) é ineficaz na implementação atual por limitação na modelagem do trocador de calor. Corrigir essa modelagem abriria acesso a distúrbios térmicos do benchmark e enriqueceria os cenários de teste do supervisor.

\end{enumerate}
